\section{Worked Example 3: Chatbot Arena 为什么用 pairwise}

Chatbot Arena 不直接让模型做标准题，而是让用户在两个匿名回复之间做选择。这样做的好处是，用户不需要把好坏翻译成一个抽象分数，只要回答“哪个更好”。

\begin{figure}[H]
\centering
\begin{tikzpicture}[node distance=1.0cm, every node/.style={font=\small, align=center}]
\node[draw, rounded corners, fill=yellow!10, minimum width=3.2cm, minimum height=0.8cm] (p) {prompt};
\node[draw, rounded corners, fill=yellow!20, above right=0.7cm and 1.4cm of p, minimum width=2.8cm, minimum height=0.8cm] (a) {model A};
\node[draw, rounded corners, fill=yellow!20, below right=0.7cm and 1.4cm of p, minimum width=2.8cm, minimum height=0.8cm] (b) {model B};
\node[draw, rounded corners, fill=orange!24, right=1.8cm of a, minimum width=3.0cm, minimum height=1.0cm] (j) {human judge};
\draw[-{Latex[length=3mm]}, thick] (p) -- (a);
\draw[-{Latex[length=3mm]}, thick] (p) -- (b);
\draw[-{Latex[length=3mm]}, thick] (a) -- (j);
\draw[-{Latex[length=3mm]}, thick] (b) -- (j);
\end{tikzpicture}
\caption{Chatbot Arena 的核心：pairwise comparison 比绝对打分更自然}
\end{figure}

它的一个关键优势是，只需要比较两个答案，不需要人类定义“这次到底是 7 分还是 8 分”。但它也有明显问题：
\begin{itemize}
    \item 输入分布是活流量，不稳定。
    \item 用户群体不是均匀抽样。
    \item ELO 是相对排名，不等于绝对能力。
\end{itemize}

\begin{knowledgebox}{为什么 pairwise 很常见}
pairwise 比绝对评分更容易做出一致判断，也更接近真实用户的“我更喜欢 A 还是 B”。缺点是你只能得到相对顺序，解释空间也更依赖流量和对局采样。
\end{knowledgebox}

